\begin{Remark}\label{Remark1}
In \cite{ZGBJPhysA}, the central element $C^{\Lambda_0}$ is defined, where $\Lambda_0$ is the highest weight of a finite-dimensional irreducible module $V(\Lambda_0)$. The construction there is similar to the one here, with the major difference being the use of explicit universal $R$-matrices in place of fused $R$-matrices. Although it is not necessarily simple to check directly that $C_m$ equals (up to a~constant) $C^{\Lambda_0}$ for $V(\Lambda_0)=P_mV^{\otimes m}$, it is straightforward to check that their eigenvalues are the same.

If $V(\Lambda_0)$ has distinct weights $\lambda_1,\dots,\lambda_r$ with multiplicities $d_1,\dots,d_r$, then the eigenvalue of~$C^{\Lambda_0}$ on an irreducible module with highest weight $\Lambda$ is given by
\[
\sum_{k=1}^r d_k q^{2(\Lambda + \rho,\lambda_k)}.
\]
Here, $\rho$ is half the sum of the positive roots, and $(\cdot,\cdot)$ is the usual invariant bilinear form on $\mathfrak{h}^*$. Now take $\Lambda_0$ to be the highest weight of $P_mV^{\otimes m}$; then the distinct weights are elements of $\mathcal{B}^{(N)}_m$ with multiplicity $1$. Therefore the eigenvalue is
\[
\sum_{\mu \in \mathcal{B}^{(N)}_m} q^{(2\rho,\mu)}q^{2(\Lambda,\mu)}.
\]

If $C_m$ acts on the same module $V(\Lambda)$, then its eigenvalue can be found by evaluating on the lowest weight vector, because then only the diagonal terms $(\ii=\jj)$ have a nonzero contribution. So the eigenvalue is
\[
\big(q-q^{-1}\big)^{-2m}q^{-Nm}\sum_{\ii \in \mathcal{W}_m} q^{ 2i_1 + \dots + 2i_m} q^{(2\mu^{(N)}(\ii),\Lambda)}.
\]
By Lemma \ref{Leem}, this is $\big(q-q^{-1}\big)^{-2m}q^{-Nm}$ times the eigenvalue of $C^{\Lambda_0}$.
\end{Remark}
